\documentclass[10pt,a4paper]{amsart}
\usepackage[margin=19mm]{geometry}
\usepackage{amsmath,amssymb,mathtools}
\usepackage[scr=rsfs,cal=euler]{mathalfa}
\usepackage{xfrac}
\usepackage[unicode,hidelinks,bookmarks=false]{hyperref}
\newcommand{\ie}{i.e.\ }
\newcommand{\cf}{cf.\ }
\newcommand{\resp}{respectively\ }
\newcommand{\Subsection}{\subsection}
\providecommand{\nobreakdash}{\nobreak\hbox{-}}
\setlength{\emergencystretch}{2em}
\multlinegap0pt
\usepackage{amssymb}
\usepackage{xcolor}
\newtheorem{corollary}{Corollary}
\newtheorem{lemma}{Lemma}
\newtheorem{theorem}{Theorem}
\title{New orientable sequences}
\author{Chris J Mitchell, Peter R Wild}
\date{Source: arXiv:2507.02526v1 (CC BY 4.0)}
\begin{document}
\maketitle

\noindent\emph{Source and licence.} New orientable sequences. Version arXiv:2507.02526v1 (CC BY 4.0), distributed by arXiv under the Creative Commons Attribution 4.0 International licence, http://creativecommons.org/licenses/by/4.0/, as recorded in the arXiv licence metadata for this exact version and shown on the abstract page https://arxiv.org/abs/2507.02526v1. Full licence notice: \texttt{licenses/arxiv-2507.02526-CC-BY-4.0.txt}.\par\smallskip

\noindent\textit{Editorial scope.} Counted unit: the article's theorem theorem:OS\_bound\_4 (source0.tex lines 649--677). The article's complete original proof of it is delivered with it (source0.tex lines 679--780, one \texttt{proof} environment of 102 lines). No mathematics is written, repaired or re-derived: the statement and the proof are the article's own lines, in the article's own notation, and no retained line is substituted or re-flowed.  Retained uncounted as the source-local context the proof needs: lines 333--349; lines 417--433; lines 455--477; lines 531--547; lines 612--627.  Nothing was omitted from any retained span, so the package carries no omission record.  No retained line contains a citation, so the wrapper carries no bibliography and no bibliography entry.  No retained line contains a figure environment, an image-inclusion command or drawing code, so no image is delivered and the package declares no asset dependency.  Editorial formatting only, in addition: an AMS \texttt{amsart} preamble that restates the article's own declarations and macros used by the retained lines, namely the environments \texttt{corollary}, \texttt{lemma}, \texttt{theorem} on separate counters, together with the standard packages those lines need.  The retained environments print in this document as Corollary 1, Lemma 1, Theorem 1 in order of appearance, whereas the article prints the same items with the numbering of its own full document.  The complete native source of the article as distributed in the arXiv e-print for this exact version is bundled unchanged as \texttt{sources/180948/source0.tex}.
\begin{corollary}  \label{corollary:excluded_tuples_by_unequal_degree}
Suppose $k\geq2$ and $n\geq4$ and $S$ is an $\mathcal{OS}_k(n)$. Then
\begin{itemize}
\item[i)] if $n$ is even, for every vertex in $B^*_k(n-1)$ corresponding to a non-uniform
    left-semi-symmetric $(n-1)$-tuple $(a_0,a_1,\ldots,a_{n-2})$, there is an edge
    $(a_0,a_1,\ldots,a_{n-2},x)$ in $B^*_k(n-1)$, for some $x$, that is not in $B(S,n)$;
\item[ii)] if $n$ is even, for every vertex in $B^*_k(n-1)$ corresponding to a non-uniform
    right-semi-symmetric $(n-1)$-tuple $(a_0,a_1,\ldots,a_{n-2})$, there is an edge
    $(y,a_0,a_1,\ldots,a_{n-2})$ in $B^*_k(n-1)$, for some $y$, that is not in $B(S,n)$;
\item[iii)] if $n$ is odd, for every vertex in $B^*_k(n-1)$ corresponding to a non-uniform
    non-alternating left-semi-symmetric $(n-1)$-tuple $(a_0,a_1,\ldots,a_{n-2})$, there is an
    edge $(a_0,a_1,\ldots,a_{n-2},x)$ in $B^*_k(n-1)$, for some $x$, that is not in $B(S,n)$;
\item[iv)] if $n$ is odd, for every vertex in $B^*_k(n-1)$ corresponding to a non-uniform
    non-alternating right-semi-symmetric $(n-1)$-tuple $(a_0,a_1,\ldots,a_{n-2})$, there is an
    edge $(y,a_0,a_1,\ldots,a_{n-2})$ in $B^*_k(n-1)$, for some $y$, that is not in $B(S,n)$.
\end{itemize}
\end{corollary}

\begin{corollary}  \label{corollary:excluded-tuples-by-parity}
Suppose $k\geq2$ and $n\geq2$ and $S$ is an $\mathcal{OS}_k(n)$. Then
\begin{itemize}
\item[i)] if $k$ is even, for every vertex in $B^*_k(n-1)$ corresponding to a uniform
    $(n-1)$-tuple $(a,a,...,a)$, there is an edge $(a,a,...,a,x)$ in $B^*_k(n-1)$, for some
    $x$, that is not in $B(S,n)$.
\item[ii)] if $k$ is even, for every vertex in $B^*_k(n-1)$ corresponding to a uniform
    $(n-1)$-tuple $(a,a,...,a)$, there is an edge $(y,a,a,...,a)$ in $B^*_k(n-1)$, for some
    $y$, that is not in $B(S,n)$.
\item[iii)] if $k$ is odd, for every vertex in $B^*_k(n-1)$ corresponding to a symmetric and
    non-uniform $(n-1)$-tuple $(a_0,a_1,\ldots,a_{n-2})$, there is an edge
    $(a_0,a_1,\ldots,a_{n-2},x)$ in $B^*_k(n-1)$, for some $x$, that is not in $B(S,n)$.
\item[iv)] if $k$ is odd, for every vertex in $B^*_k(n-1)$ corresponding to a symmetric and
    non-uniform $(n-1)$-tuple $(a_0,a_1,\ldots,a_{n-2})$, there is an edge
    $(y,a_0,a_1,\ldots,a_{n-2})$ in $B^*_k(n-1)$, for some $y$, that is not in $B(S,n)$.
\end{itemize}
\end{corollary}

\begin{lemma}  \label{lemma:constraint-interactions}
Suppose $k\geq2$.  Suppose $\mathbf{a}=(a_0,a_1,\dots,a_{n-2})$ and
$\mathbf{b}=(a_1,a_2,\dots,a_{n-1})$ are $k$-ary $(n-1)$-tuples, connected by the edge
$(a_0,a_1,\dots,a_{n-1})$.
\begin{itemize}
\item[i)] If $n\geq3$, $\mathbf{a}$ is symmetric and $\mathbf{b}$ is symmetric then either
    $\mathbf{a}$ and $\mathbf{b}$ are both uniform or $n$ is even and $\mathbf{a}$ and
    $\mathbf{b}$ are both alternating.
\item[ii)] If $n\geq5$, $\mathbf{a}$ is symmetric and $\mathbf{b}$ is right-semi-symmetric then
    there exist $c_j$, $0\leq j\leq2$, such that $a_{3i+j}=c_j$, for every $i$ and $0\leq j\leq
    2$.  Moreover if $n\equiv0\pmod3$ then $c_0=c_1$, if $n\equiv1\pmod3$ then $c_0=c_2$, and
    if $n\equiv2\pmod3$ then $c_1=c_2$.
\item[iii)] If $n\geq5$, $\mathbf{a}$ is left-semi-symmetric and $\mathbf{b}$ is symmetric then
    there exist $c_j$, $0\leq j\leq2$, such that $a_{3i+j}=c_j$, for every $i$ and $0\leq j\leq
    2$.  Moreover if $n\equiv0\pmod3$ then $c_1=c_2$, if $n\equiv1\pmod3$ then $c_0=c_1$, and
    if $n\equiv2\pmod3$ then $c_0=c_2$.
\item[iv)] If $n\geq5$, $\mathbf{a}$ is left-semi-symmetric and $\mathbf{b}$ is
    right-semi-symmetric then there exist $c_j$, $0\leq j\leq3$, such that $a_{4i+j}=c_j$, for
    every $i$ and $0\leq j\leq 3$. Moreover if $n\equiv0\pmod4$ then $c_0=c_1$ and $c_2=c_3$,
    if $n\equiv1\pmod4$ then $c_0=c_2$, if $n\equiv2\pmod4$ then $c_0=c_3$ and $c_1=c_2$, and
    if $n\equiv3\pmod4$ then $c_1=c_3$.
\end{itemize}
\end{lemma}

\begin{lemma}  \label{lemma:tuple_numbers}
Suppose $n\geq3$ and $k\geq2$.  Then
\begin{itemize}
\item[i)] the number of symmetric $k$-ary $n$-tuples is $k^{\lceil n/2\rceil}$;
\item[ii)] the number of uniform $k$-ary $n$-tuples is $k$.
\item[iii)] the number of symmetric non-uniform $k$-ary $n$-tuples is $k^{\lceil n/2\rceil}-k$;
\item[iv)] the number of asymmetric $k$-ary $n$-tuples is $k^n-k^{\lceil n/2\rceil}$;
\item[v)] the number of alternating $k$-ary $n$-tuples is $k(k-1)$;
\item[vi)]  the number of non-uniform non-alternating symmetric $k$-ary $n$-tuples is
    $k^{(n+1)/2}-k^2$ if $n$ is odd, and $k^{n/2}-k$ if $n$ is even;
\item[vii)] the number of left-semi-symmetric $k$-ary $n$-tuples is $k^{\lceil(n+1)/2\rceil}$;
\item[viii)] the number of non-uniform left-semi-symmetric $k$-ary $n$-tuples is
    $k^{\lceil(n+1)/2\rceil}-k$;
\item[ix)] the number of non-uniform non-alternating left-semi-symmetric $k$-ary $n$-tuples is
    $k^{(n+2)/2}-k^2$ if $n$ is even, and $k^{(n+1)/2}-k$ if $n$ is odd.
\end{itemize}
\end{lemma}

\begin{lemma}  \label{lemma:bounds-high-level}
If $k\geq2$ and $S$ is an $\mathcal{OS}_k(n)$ then:
\[ |B(S,n)|\leq \begin{cases}
k^2-k - |P_{\text{out}}|,
 & \text{if $n=2$} \\
k^3-k^2 - |P_{\text{out}}|-|P_{\text{in}}| + |P_{\text{out}}\cap P_{\text{in}}|,
 & \text{if $n=3$} \\
k^4-k^2 - |U_{\text{out}}|-|P_{\text{out}}|,
 & \text{if $n=4$} \\
k^n-k^{\lceil n/2\rceil} - |U_{\text{out}}|-|U_{\text{in}}|-|P_{\text{out}}|-|P_{\text{in}}| \\
~~~~+~|U_{\text{out}}\cap P_{\text{in}}|+|P_{\text{out}}\cap U_{\text{in}}| \\
~~~~+~|U_{\text{out}}\cap U_{\text{in}}|+|P_{\text{out}}\cap P_{\text{in}}|,
 & \text{if $n>4$.}
\end{cases} \]
where $|X\cap Y|$ denotes the maximum possible cardinality for such a set.
\end{lemma}

\begin{theorem} \label{theorem:OS_bound_4}
Suppose $k\geq2$ and $m$ is the period of an $\mathcal{OS}_k(n)$.
\begin{itemize}
\item[i)] If $n=2$ then:
\[ m \leq \begin{cases}
\frac{k^2-k}{2}, & \text{if $k$ is odd} \\
\frac{k^2-2k}{2}, & \text{if $k$ is even}
\end{cases} \]
\item[ii)] If $n=3$ then:
\[ m \leq \begin{cases}
\frac{k^3-k^2}{2}, & \text{if $k$ is odd} \\
\frac{k^3-k^2-2k}{2}, & \text{if $k$ is even}
\end{cases} \]
\item[iii)] If $n=4$ then:
\[ m \leq \begin{cases}
\frac{k^4-3k^2+2k}{2}, & \text{if $k$ is odd} \\
\frac{k^4-2k^2}{2}, & \text{if $k$ is even}
\end{cases}\]
\item[iv)] If $n>4$ then:
\[ m \leq \begin{cases}
\frac{k^n-3k^{(n+1)/2}-2k^{(n-1)/2}+k^3+3k^2}{2}, & \text{if $n$ and $k$ are odd} \\
\frac{k^n-3k^{(n+1)/2}+k^3+k^2-2k}{2}, & \text{if $n$ is odd and $k$ is even} \\
\frac{k^n-5k^{n/2}+4k^2}{2}, & \text{if $n$ is even and $k$ is odd} \\
\frac{k^n-3k^{n/2}+k^2-k}{2}, & \text{if $n$ and $k$ are even}
\end{cases}\]

\end{itemize}

\end{theorem}

\begin{proof}
The proof builds on, and uses the notation of, Lemma~\ref{lemma:bounds-high-level}.
\begin{itemize}

\item[i)] {\bf $n=2$}.  If $k$ is odd then $P_{\text{out}}=\emptyset$ by
    Corollary~\ref{corollary:excluded-tuples-by-parity}(iii), since there are no symmetric
    non-uniform 1-tuples.  If $k$ is even then, by
    Corollary~\ref{corollary:excluded-tuples-by-parity}(i), $|P_{\text{out}}|=k$, since there
    are $k$ uniform 1-tuples by Lemma~\ref{lemma:tuple_numbers}(ii). The result follows from
    Lemma~\ref{lemma:bounds-high-level}.

\item[ii)] {\bf $n=3$}.  If $k$ is odd then, by
    Corollary~\ref{corollary:excluded-tuples-by-parity}(iii),(iv),
    $|P_{\text{in}}|=|P_{\text{out}}|$ equals the number of symmetric non-uniform 2-tuples,
    i.e.\ zero by Lemma~\ref{lemma:tuple_numbers}(iii).

    If $k$ is even then, by Corollary~\ref{corollary:excluded-tuples-by-parity}(i),(ii),
    $|P_{\text{in}}|=|P_{\text{out}}|$ equals the number of uniform 2-tuples, i.e.\ $k$ by
    Lemma~\ref{lemma:tuple_numbers}(ii). Moreover, $P_{\text{in}}\cap
    P_{\text{out}}=\emptyset$, since an edge can only be outgoing from a uniform ($n-1)$-tuple
    and incoming to a uniform ($n-1)$-tuple if it corresponds to a uniform $n$-tuple, and such
    an edge cannot occur in $B^*_k(n-1)$. The result follows from
    Lemma~\ref{lemma:bounds-high-level}.

\item[iii)] {\bf $n=4$}. By Corollary~\ref{corollary:excluded_tuples_by_unequal_degree}(i),
    $|U_{\text{out}}|$ equals the number of non-uniform left-semi-symmetric $3$-tuples, i.e.\
    $k^2-k$ by Lemma~\ref{lemma:tuple_numbers}(viii).  If $k$ is odd then, by
    Corollary~\ref{corollary:excluded-tuples-by-parity} (iii), $|P_{\text{out}}|$ equals the
    number of symmetric non-uniform 3-tuples, i.e.\ $k^2-k$ by
    Lemma~\ref{lemma:tuple_numbers}(iii). If $k$ is even then, by
    Corollary~\ref{corollary:excluded-tuples-by-parity} (i), $|P_{\text{out}}|$ equals the
    number of uniform 3-tuples, i.e.\ $k$ by Lemma~\ref{lemma:tuple_numbers}(ii). The result
    follows from Lemma~\ref{lemma:bounds-high-level}.



\item[iv)a)] {\bf $n>4$; $n$ odd}.  By
    Corollary~\ref{corollary:excluded_tuples_by_unequal_degree}(iii),(iv),
    $|U_{\text{out}}|=|U_{\text{in}}|$ equals the number of non-uniform non-alternating
    left-semi-symmetric $(n-1)$-tuples, i.e.\ $k^{(n+1)/2}-k^2$ by
    Lemma~\ref{lemma:tuple_numbers}(ix).  Also, by
    Lemma~\ref{lemma:constraint-interactions}(iv), $|U_{\text{out}}\cap U_{\text{in}}|=k^3-k^2$
    since if, in the statement of the lemma, we choose $c_1\neq c_3$ if $n\equiv 1\pmod 4$, and
    $c_0\neq c_2$ if $n\equiv 3\pmod 4$, then $\mathbf{a}$ and $\mathbf{b}$ are non-uniform and
    non-alternating and there are $k^2(k-1)=k^3-k^2$ possibilities for the values $c_i$.

    If $k$ is odd then, by Corollary~\ref{corollary:excluded-tuples-by-parity} (iii),(iv),
    $|P_{\text{out}}|=|P_{\text{in}}|$ equals the number of symmetric non-uniform
    $(n-1)$-tuples, i.e.\ $k^{(n-1)/2}-k$ by Lemma~\ref{lemma:tuple_numbers}(iii).\newline By
    Lemma~\ref{lemma:constraint-interactions}(ii),(iii), $|U_{\text{out}}\cap
    P_{\text{in}}|=|P_{\text{out}}\cap U_{\text{in}}|=k(k-1)$ (removing the case where the
    $c_i$ are all equal).
    \newline By Lemma~\ref{lemma:constraint-interactions}(i), $P_{\text{out}}\cap
    P_{\text{in}}=\emptyset$.

    If $k$ is even then, by Corollary~\ref{corollary:excluded-tuples-by-parity}(i),(ii),
    $|P_{\text{out}}|=|P_{\text{in}}|$ equals the number of uniform $(n-1)$-tuples, i.e.\ $k$
    by Lemma~\ref{lemma:tuple_numbers}(ii).
    \newline By Corollary~\ref{corollary:excluded-tuples-by-parity}(i),(ii), $P_{\text{out}}$
    and $P_{\text{in}}$ only contain edges out-going/in-going from/to uniform $(n-1)$-tuples,
    and by Corollary~\ref{corollary:excluded_tuples_by_unequal_degree}(iii),(iv)
    $U_{\text{out}}$ and $U_{\text{in}}$ only contain edges that are out-going/in-going from/to
    non-uniform $(n-1)$-tuples, and hence $|U_{\text{out}}\cap
    P_{\text{in}}|=|P_{\text{out}}\cap U_{\text{in}}|=\emptyset$.
    \newline By Corollary~\ref{corollary:excluded-tuples-by-parity}(i),(ii), $P_{\text{out}}$
    and $P_{\text{in}}$ only contain edges out-going/in-going from/to uniform $(n-1)$-tuples,
    and hence an edge in $P_{\text{out}}\cap P_{\text{in}}$ must be uniform, i.e.\
    $P_{\text{out}}\cap P_{\text{in}}=\emptyset$.

    The result follows from Lemma~\ref{lemma:bounds-high-level}.

\item[iv)b)] {\bf $n>4$; $n$ even}. By
    Corollary~\ref{corollary:excluded_tuples_by_unequal_degree}(i),(ii),
    $|U_{\text{out}}|=|U_{\text{in}}|$ equals the number of non-uniform left-semi-symmetric
    $(n-1)$-tuples, i.e.\ $k^{n/2}-k$ by Lemma~\ref{lemma:tuple_numbers}(viii). Additionally,
    by Lemma~\ref{lemma:constraint-interactions}(iv), $|U_{\text{out}}\cap
    U_{\text{in}}|=k(k-1)$, by choosing $c_0$ and $c_2$ to be distinct.

    If $k$ is odd then, by Corollary~\ref{corollary:excluded-tuples-by-parity}(iii),(iv),
    $|P_{\text{out}}|=|P_{\text{in}}|$ equals the number of symmetric non-uniform
    $(n-1)$-tuples, i.e.\ $k^{n/2}-k$ by Lemma~\ref{lemma:tuple_numbers}(iii).
    \newline By Lemma~\ref{lemma:constraint-interactions}(ii),(iii), $|U_{\text{out}}\cap
    P_{\text{in}}|=|P_{\text{out}}\cap U_{\text{in}}|=k(k-1)$, by ensuring that $c_0$, $c_1$
    and $c_2$ are not all equal.
    \newline By Lemma~\ref{lemma:constraint-interactions}(i), $|P_{\text{out}}\cap P_{\text{in}}|=k(k-1)$,
    i.e.\ the number of non-uniform alternating $(n-1)$-tuples.

    If $k$ is even then, by Corollary~\ref{corollary:excluded-tuples-by-parity}(i),(ii),
    $|P_{\text{out}}=|P_{\text{in}}|$ equals the number of uniform $(n-1)$-tuples, i.e.\ $k$ by
    Lemma~\ref{lemma:tuple_numbers}(i).
    \newline By Corollary~\ref{corollary:excluded-tuples-by-parity}(i),(ii), $P_{\text{out}}$
    and $P_{\text{in}}$ only contain edges out-going/in-going from/to uniform $(n-1)$-tuples,
    and by Corollary~\ref{corollary:excluded_tuples_by_unequal_degree}(i),(ii) $U_{\text{out}}$
    and $U_{\text{in}}$ only contain edges that are out-going/in-going from/to non-uniform
    $(n-1)$-tuples, and hence $|U_{\text{out}}\cap P_{\text{in}}|=|P_{\text{out}}\cap
    U_{\text{in}}|=\emptyset$.
    \newline Finally, it is immediate that $P_{\text{out}}\cap P_{\text{in}}=\emptyset$.

    The result follows from Lemma~\ref{lemma:bounds-high-level}.

\end{itemize}
\end{proof}

\end{document}
